% !TeX root = ../main.tex
\begin{frame}{02 研究现状——支架构型}
  \begin{SlideCanvas}
    \MiddleTimeAxis
    \MiddleTimeNode{155}{2018}{up}
    \MiddleTimeNode{370}{2020}{down}
    \MiddleTimeNode{585}{2025}{up}
    \MiddleTimeNode{800}{2026}{down}
    \PointAt{55}{85}{400}{Wu等：反四手性}
    \CenteredPicAt{145}{180}{180}{112}{assets/wu2018_fig1.jpg}
    \DetailAt{250}{126}{205}{圆形／椭圆节点和层级构型调节整体变形。}
    \PointAt{505}{85}{400}{Hasanzadeh等：混合单胞}
    \CenteredPicAt{595}{180}{180}{112}{assets/has2025_fig2a.png}
    \DetailAt{700}{126}{205}{组合箭头与缺肋几何，形成混合负泊松比单胞并构建管状支架。}
    \PointAt{55}{333}{400}{Ruan等：周向约束}
    \CenteredPicAt{145}{423}{180}{105}{assets/ruan2020_ps1.png}
    \DetailAt{250}{376}{205}{引入周向约束杆调节变形耦合，并评价狭窄血管中的部署表现。}
    \PointAt{505}{333}{400}{Zhang等：花瓣手性}
    \CenteredPicAt{595}{423}{180}{105}{assets/zhang2026_fig17.jpg}
    \DetailAt{700}{376}{205}{设计花瓣形手性单胞，结合参数反设计与球囊／血管部署仿真。}
    \Sources{\RefWu，图1；\RefRuan，图1中PS1。\\\RefHas，图2(a)；\RefZhang，图17。年份按卷期，间距示意；Zhang在线发表于2025年。}
  \end{SlideCanvas}
  \BICConclusionBand{支架构型通过\BICBandEmph{单胞形态和连接约束}调节整体变形}
  \note{四篇工作沿时间轴上下错列，每篇一张对应原图。Wu2018介绍圆形／椭圆节点与层级反四手性；Ruan2020介绍约束杆对周向耦合及部署的影响。Hasanzadeh2025结合arrowhead与missing-rib单胞，使用原文图2(a)的完整构型生成面板，包含平面到圆管的过程；并非新的计算或本课题组成果。Zhang2026展示花瓣手性、参数反设计和部署模型，其神经网络／差分进化不混称材料分布拓扑优化。文献介绍依据公开摘要、预览和原图；不推断临床效果。年份按卷期，节点间距不表示实际时间长度。}
\end{frame}
